\specsection{Measures on general topological spaces and tight convergence}

The study of the connections between topology and measure theory has above all been conceived of as the study of regularity properties of measures, and in particular that of `outer' regularity and of `inner' regularity; \(^{(14)}\) inner regularity is equivalent to outer regularity on a locally compact space countable at infinity. The construction that Lebesgue gives to the measure of subsets of the line highlights these two kinds of regularity, and the outer regularity property of measures on a Polish space seems to have had public notoriety around 1935. But it was not until 1940, in an article whose diffusion was retarded by the war, that A. D. Alexandroff (XIX) gave prominence to the role of inner regularity and showed that it is possessed by the measures on a Polish space; this result was rediscovered later by Prokhorov (XIII) and is often erroneously attributed to this author. It was not perceived until very recently that this property extends to Souslin spaces; because of this, the importance of these spaces has increased greatly, even more so as it was realized that their theory could be worked out without any metrizability hypothesis, and that the quasi-totality of the function spaces were Souslin (most often, even Lusin). \(^{(15)}\) These are the reasons that moved us to place the accent on inner regular measures in this chapter.

The definition of a mode of convergence (vague or tight) for measures is most conveniently done by putting the space of measures into duality with a space of continuous functions. Generalizing an old result of F. Riesz, A. A. Markoff established in 1938 a one-to-one correspondence between the positive functionals on \(\mathcal{C}(\mathrm{X})\) and the regular measures on a compact space X. In the work (XIX) already cited, A. D. Alexandroff extends these results to the case of a completely regular space: he introduces a hierarchy in the set of positive linear forms on the space \(\mathcal{C}^b (\mathrm{X})\) of bounded continuous functions on a completely regular space \(\mathrm{X},^{(16)}\) he defines tight convergence of bounded measures and proves among others the following two theorems:

\begin{enumerate}
\def\labelenumi{\alph{enumi})}
\item
  if X is Polish, the set of linear forms on \(\mathcal{C}^{b}(X)\) corresponding to the measures is closed for the weak convergence of sequences;
\item
  if a sequence of bounded measures has a tight limit, `no mass escapes at infinity' (this is a weak form of the converse of Prokhorov's theorem on tight convergence).
\end{enumerate}

From this abundance of concepts and theorems, Prokhorov was able to extract the results important for the theory of stochastic processes, and to present them in a simple and striking form. In his great work of 1956 already cited (XIII), a large part is devoted to bounded positive measures on a Polish space; generalizing a construction of Lévy, he defines a metric on the set of positive measures of mass 1 that makes it a Polish space, and then establishes an important compactness criterion for tight convergence (cf.~§5, No.~5, Th. 1). Independently of Prokhorov, Le Cam (XX) obtained a number of compactness results for the tight convergence of measures; he makes no metrizability hypothesis on the spaces he considers, and his results reduce to earlier theorems of Dieudonné in the locally compact case.

